\chapter{Fret, météo et autres marchés de niche}\label{ch:m3:freight-weather-and-other-corners}

Une entreprise de distribution de gaz à Chicago gagne son argent en hiver. Un hiver froid remplit ses canalisations et ses
comptes~; un hiver doux la laisse avec du gaz acheté pour des clients qui n'en avaient pas besoin.
Pendant l'hiver 2011/12, l'aéroport de la ville a enregistré 3\,932 \omterm{def:m3:freight-weather-and-other-corners:dd}{degrés-jours de chauffage}, le plus faible en
trente-six hivers~; deux ans plus tard, l'aéroport en a enregistré 5\,999, le plus élevé. Le distributeur peut acheter un
contrat qui le rémunère pour chaque degré-jour dont l'hiver reste en deçà
d'un prix d'exercice, réglé sur les températures relevées à cet aéroport. Le fret et la météo sont des marchés portant sur
des choses qui ne sont pas du tout des matières premières~: l'usage d'un navire pour une traversée, la température
d'une saison. Ce chapitre expose comment les navires sont affrétés et comment le fret est couvert, comment
la météo est transformée en indice négocié et valorisée, et ce qu'il faut pour qu'un tel marché de niche
devienne un marché.

\section{Fret~: affrètements et routes maritimes}

Les matières premières en vrac voyagent à bord de navires affrétés à cet effet. L'affrètement prend deux formes.

\begin{definition}[Affrètement au voyage, affrètement à temps]\label{def:m3:freight-weather-and-other-corners:charter}
Un \emph{affrètement au voyage}\index{affrètement au voyage} loue un navire pour transporter une cargaison entre des ports
désignés, moyennant un fret par tonne~; l'armateur prend en charge les frais de voyage (carburant, frais portuaires). Un
\emph{affrètement à temps}\index{affrètement à temps} loue un navire pour une période, moyennant un taux journalier de location~;
l'affréteur dirige ses voyages et en paie les coûts.
\end{definition}

\begin{definition}[Équivalent affrètement à temps]\label{def:m3:freight-weather-and-other-corners:tce}
L'\emph{équivalent affrètement à temps}\index{equivalent affrètement à temps@équivalent affrètement à temps} (TCE) d'un voyage est sa recette de fret
diminuée de ses frais de voyage, divisée par sa durée en jours~: le taux journalier de location auquel un
\omterm{def:m3:freight-weather-and-other-corners:charter}{affrètement à temps} rapporterait à l'armateur le même revenu.
\end{definition}

\begin{example}[Du voyage au taux journalier]\label{ex:m3:freight-weather-and-other-corners:tce}
Un vraquier transporte 170\,000 tonnes de minerai de fer à 12.00~\$ la tonne, soit un fret de 2.04 millions de dollars.
Le carburant coûte 600\,000~\$ et les frais portuaires 150\,000~\$, et le voyage aller-retour dure 45 jours. Son TCE est
$(2\,040\,000 - 750\,000)/45 = \$28\,667$ par jour.
\end{example}

\begin{omfigure}
\centering
\begin{tikzpicture}[font=\small,
  pty/.style={draw,thick,rounded corners=2pt,align=center,minimum width=2.6cm,minimum height=1.1cm},
  arr/.style={-{Stealth[length=2.2mm]},thick}]
\node[pty,draw=omDef,fill=omDef!8] (own) at (0,0) {armateur};
\node[pty,draw=omThm,fill=omThm!8] (chr) at (9.6,0) {affréteur};
\draw[arr] ([yshift=4pt]chr.west) -- node[above=2pt,font=\footnotesize] {affrètement au voyage~: \$ par tonne} ([yshift=4pt]own.east);
\draw[arr] ([yshift=-4pt]chr.west) -- node[below=2pt,font=\footnotesize] {affrètement à temps~: \$ par jour} ([yshift=-4pt]own.east);
\node[font=\footnotesize,align=center,anchor=north] at (0,-0.9) {prend en charge carburant et ports\\sous un affrètement au voyage};
\node[font=\footnotesize,align=center,anchor=north] at (9.6,-0.9) {prend en charge carburant et ports\\sous un affrètement à temps};
\end{tikzpicture}

\omcaption{Les deux façons d'affréter un navire, et celui qui supporte les frais de voyage.
L'\omterm{def:m3:freight-weather-and-other-corners:tce}{équivalent affrètement à temps} permet de passer de l'une à l'autre. Schéma indicatif.}\label{fig:m3:freight-weather-and-other-corners:charters}
\end{omfigure}

\begin{definition}[Indice de route]\label{def:m3:freight-weather-and-other-corners:route}
Un \emph{indice de route}\index{indice de route} est une évaluation quotidienne, réalisée par un panel de courtiers maritimes
selon les règles du fournisseur d'indice, du taux d'un navire standard sur une route standard, en dollars par tonne ou
par jour~; les indices de plusieurs routes sont moyennés en un panier (pour les vraquiers Capesize, une
moyenne pondérée de cinq routes d'\omterm{def:m3:freight-weather-and-other-corners:charter}{affrètement à temps}).
\end{definition}

\begin{dated}{2026-09}{Le fournisseur des indices de fret}\label{dat:m3:freight-weather-and-other-corners:baltic}
Le Baltic Exchange, qui publie les principaux indices de fret sec et humide, a été racheté par Singapore
Exchange~; le rachat a été achevé en novembre 2016.
\end{dated}

\section{Accords de fret à terme}

\begin{definition}[Forward freight agreement (FFA)]\label{def:m3:freight-weather-and-other-corners:ffa}
Un \emph{forward freight agreement}\index{forward freight agreement} (FFA) est un contrat à terme réglé en espèces sur
un \omterm{def:m3:freight-weather-and-other-corners:route}{indice de route} pour une période future, en général un mois~: à l'échéance, l'acheteur reçoit, par jour (ou
par tonne) de la quantité convenue, la moyenne de l'indice sur la période moins le taux fixe convenu.
La plupart font l'objet d'une compensation.
\end{definition}

Un armateur qui prévoit de louer son navire dans six mois vend des FFA pour verrouiller son taux de location~; un affréteur ou
une \omterm{def:m3:physical-commodity-markets:house}{maison de négoce} qui aura besoin de navires en achète. Comme les FFA se règlent sur une moyenne mensuelle, ils couvrent
la moyenne du taux du mois, et non le taux du jour où un navire est fixé~: cet écart constitue le \omterm{def:m3:physical-commodity-markets:basisrisk}{risque de base}
(\cref{ch:m3:physical-commodity-markets}).

\begin{omfigure}
\begin{tikzpicture}
\begin{axis}[width=11cm,height=4.4cm,
  xlabel={jour de l'indice dans le mois},ylabel={milliers de \$ par jour},
  tick label style={font=\small},label style={font=\small},
  xmin=0,xmax=26,ymin=20,ymax=27,grid=major,grid style={gray!25},
  legend columns=3,legend style={font=\footnotesize,at={(0.5,-0.33)},anchor=north,draw=none,fill=none,/tikz/every even column/.append style={column sep=8pt}},
  table/col sep=comma]
\addplot[omDef,very thick,mark=*,mark size=1pt] table[x=day,y=index]{figdata/markets-3/11-freight-weather-and-other-corners/ffa.csv};
\addplot[omThm,thick,dashed,domain=0:26] {23.4};
\addplot[omProp,thick,dotted,domain=0:26] {21.0};
\legend{indice quotidien,moyenne mensuelle (règlement),taux fixe du FFA}
\end{axis}
\end{tikzpicture}

\omcaption{Mois de règlement d'un FFA à \omterm{def:m3:freight-weather-and-other-corners:charter}{affrètement à temps}, à titre illustratif~: l'indice passe de 21\,000~\$ à
25\,800~\$ par jour sur 25 jours de l'indice, la moyenne s'établit à 23\,400~\$, et l'acheteur de 30 jours à 21\,000~\$
reçoit 72\,000~\$. Données~: tutoriel du chapitre.}\label{fig:m3:freight-weather-and-other-corners:ffa}
\end{omfigure}

\begin{dated}{2026-09}{Dérivés du fret}\label{dat:m3:freight-weather-and-other-corners:ffa}
Les contrats à terme sur fret Capesize et les FFA compensés sur Singapore Exchange se règlent en espèces sur la
moyenne arithmétique des évaluations quotidiennes du Baltic Exchange pour le panier de cinq routes d'\omterm{def:m3:freight-weather-and-other-corners:charter}{affrètement à temps}, sur
le mois du contrat~; les mêmes indices sont compensés comme contrats à terme par d'autres bourses.
\end{dated}

\section{Météo~: degrés-jours et dérivés climatiques}

La demande d'énergie suit la température, et celle-ci peut être mesurée sans contestation à une
station désignée. L'indice de la profession transforme les températures d'une saison en un seul nombre.

\begin{definition}[Degré-jour de chauffage, degré-jour de climatisation]\label{def:m3:freight-weather-and-other-corners:dd}
Avec $\theta_i$ la moyenne des températures maximale et minimale d'une journée, en degrés Fahrenheit, les
\emph{degrés-jours de chauffage}\index{degrés-jours de chauffage} de cette journée sont $\mathrm{HDD}_i = \max(0, 65 - \theta_i)$
et ses \emph{degrés-jours de climatisation}\index{degrés-jours de climatisation} $\mathrm{CDD}_i = \max(0, \theta_i - 65)$~;
l'indice d'une période est la somme de ses valeurs journalières.
\end{definition}

\begin{definition}[Dérivé climatique]\label{def:m3:freight-weather-and-other-corners:wd}
Un \emph{dérivé climatique}\index{dérivé climatique} est un contrat dont le paiement est fonction
d'un indice météorologique relevé à une station désignée sur une période~: un swap qui verse un tick par degré-jour au-dessus ou au-dessous
d'un prix d'exercice, ou une option qui ne paie que d'un seul côté, généralement avec un cap.
\end{definition}

\begin{dated}{2026-09}{Dérivés climatiques cotés en bourse}\label{dat:m3:freight-weather-and-other-corners:cme}
Les \omterm{def:m3:freight-weather-and-other-corners:wd}{dérivés climatiques} se négocient de gré à gré depuis 1997~; la première transaction autonome est décrite
comme un swap HDD entre Koch Energy et Enron pour un hiver de Milwaukee. Le Chicago Mercantile Exchange a admis à la cote des contrats à terme HDD
en septembre 1999, puis des contrats à terme CDD en janvier 2000, sur des indices à base 65°F. Son règlement de marché énumère treize
villes américaines, chacune étant relevée à une station météorologique désignée sur un aéroport (Chicago à O'Hare), additionne les degrés-jours d'un mois civil, et
fixe la valeur du contrat à 20~USD par point d'indice, avec un tick d'un point.
\end{dated}

\begin{omfigure}
\begin{tikzpicture}
\begin{axis}[width=11cm,height=4.8cm,ybar,bar width=4pt,
  xlabel={hiver (année de janvier)},ylabel={HDD, novembre--mars},
  tick label style={font=\small},label style={font=\small},
  xmin=1990,xmax=2027,ymin=3000,ymax=6500,grid=major,grid style={gray!25},
  xtick={1991,1996,2001,2006,2011,2016,2021,2026},xticklabel style={/pgf/number format/1000 sep={}},
  table/col sep=comma]
\addplot[fill=omDef!55,draw=omDef] table[x=winter,y=hdd]{figdata/markets-3/11-freight-weather-and-other-corners/hdd.csv};
\draw[omThm,thick,dashed] (axis cs:1990,4911.8) -- (axis cs:2027,4911.8);
\node[omThm,font=\footnotesize,anchor=south west] at (axis cs:1990.5,6000) {moyenne 4\,912};
\end{axis}
\end{tikzpicture}

\omcaption{\omterm{def:m3:freight-weather-and-other-corners:dd}{Degrés-jours de chauffage} à Chicago O'Hare, de novembre à mars, hivers 1990/91 à 2025/26 (moyenne
journalière des températures maximale et minimale, base 65°F). L'hiver le plus doux fut celui de 2011/12, le plus froid celui de 2013/14~; la
tendance ajustée est de $-125$ HDD par décennie. Données~: NOAA GHCN-Daily, station USW00094846.}\label{fig:m3:freight-weather-and-other-corners:hdd}
\end{omfigure}

\section{Valorisation de la météo~: burn analysis}

Un indice de degrés-jours n'a pas d'actif sous-jacent pour s'y couvrir, ni de prix à terme à répliquer~: la
météo ne se stocke ni ne s'échange. La valorisation part donc de l'historique.

\begin{definition}[Burn analysis (analyse historique)]\label{def:m3:freight-weather-and-other-corners:burn}
La \emph{burn analysis}\index{burn analysis} valorise un contrat climatique en appliquant son paiement aux valeurs
de l'indice des saisons passées, comme si chacune s'était de nouveau produite, et en prenant le paiement moyen comme fair value,
la distribution des paiements constituant son risque.
\end{definition}

\begin{method}[Burn analysis avec tendance]\label{met:m3:freight-weather-and-other-corners:burn}
(i) Calculer l'indice pour chacune des $n$ dernières saisons complètes à la station du contrat.
(ii) Ajuster une tendance linéaire et ramener chaque saison à la valeur de cette tendance pour la saison valorisée (les
changements de climat et de station font paraître les anciens hivers plus froids que ne le sera le suivant). (iii) Appliquer le paiement à chaque
saison ajustée. (iv) La fair value est le paiement moyen~; coter une marge au-dessus de celle-ci pour le risque, mesurée
par la fréquence et les quantiles élevés des paiements.
\end{method}

\begin{example}[L'option de vente du fournisseur d'énergie]\label{ex:m3:freight-weather-and-other-corners:put}
Sur les trente derniers hivers à O'Hare, ramenés à la tendance de l'hiver 2026/27, l'estimation burn des
HDD de la saison est de 4\,677. Une option de vente dont le prix d'exercice est inférieur de 10~\%, soit 4\,209, versant 20\,000~\$ par HDD avec un
plafond de 20~millions de dollars, aurait versé dans 20~\% des hivers ajustés~: 0.61~million de dollars en moyenne,
1.60~million de dollars dans l'hiver du 90\up{e}~centile et 9.30~million de dollars dans le plus doux. Sans correction de tendance,
l'estimation est de 4\,871 et le paiement moyen de 0.56~million de dollars.
\end{example}

\begin{omfigure}
\begin{tikzpicture}
\begin{axis}[width=11cm,height=4.4cm,ybar,bar width=5pt,
  xlabel={hiver (ramené à 2026/27)},ylabel={paiement, millions de \$},
  tick label style={font=\small},label style={font=\small},
  xmin=1996,xmax=2027,ymin=0,ymax=10,grid=major,grid style={gray!25},
  xtick={1997,2002,2007,2012,2017,2022,2026},xticklabel style={/pgf/number format/1000 sep={}},
  table/col sep=comma]
\addplot[fill=omThm!55,draw=omThm] table[x=winter,y=payout]{figdata/markets-3/11-freight-weather-and-other-corners/burn.csv};
\end{axis}
\end{tikzpicture}

\omcaption{\omterm{def:m3:freight-weather-and-other-corners:burn}{Burn analysis} de l'option de vente du fournisseur d'énergie~: le paiement qu'elle aurait reçu pour chacun des
trente derniers hivers, ramené au niveau de la tendance pour 2026/27. Six hivers donnent lieu à un paiement~; 2011/12 est celui qui paie le plus. Données~:
le tutoriel du chapitre, à partir de NOAA GHCN-Daily.}\label{fig:m3:freight-weather-and-other-corners:burn}
\end{omfigure}

La \omterm{def:m3:freight-weather-and-other-corners:burn}{burn analysis} est sommaire~: trente saisons en disent peu sur un hiver cinquantennal, et la tendance est elle-même
incertaine. Les desks l'affinent en simulant les températures quotidiennes et en intégrant la prévision saisonnière
lorsqu'elle est disponible, mais le prix reste un prix d'assurance, fixé par le capital de ceux qui
souscrivent le risque.

\section{Ce qui rend un marché négociable}

Le fret et la météo sont des niches parce que seules quelques conditions en ont fait des marchés, et les marchés des autres niches
ont échoué faute d'une seule. Les marchés durables en partagent cinq~:
\begin{enumerate}
\item \emph{Un besoin de couverture des deux côtés}~: armateurs et affréteurs, fournisseurs d'énergie et traders qui
sont perdants et gagnants selon la douceur de l'hiver.
\item \emph{Un indice que personne ne peut manipuler}~: une évaluation par panel selon des règles publiées, un thermomètre dans un
aéroport.
\item \emph{Des contrats standard}~: routes, stations, périodes et ticks qui concentrent la liquidité.
\item \emph{Compensation}~: une marge qui permet à des inconnus de négocier sans lignes de crédit.
\item \emph{Des spéculateurs qui acceptent de porter le risque}~: sans eux, les hedgers doivent trouver une contrepartie exacte,
ce qui arrive rarement.
\end{enumerate}
Un contrat qui n'a pas la deuxième condition échoue lorsque l'indice peut être influencé~; celui qui n'a pas la première ou la
cinquième ne se négocie jamais.

\section{Tutoriel~: degrés-jours et prix burn}\label{tut:m3:freight-weather-and-other-corners:burn}

\begin{tutorial}
\textbf{Objectif.} À partir des températures quotidiennes de la station, construire les totaux d'HDD par hiver, les corriger de la tendance et valoriser une
option de vente sur degrés-jours par \omterm{def:m3:freight-weather-and-other-corners:burn}{burn analysis}. \textbf{État final~:}
les \cref{fig:m3:freight-weather-and-other-corners:hdd,fig:m3:freight-weather-and-other-corners:burn} et les
chiffres de l'\cref{ex:m3:freight-weather-and-other-corners:put}.
\begin{tutsteps}
\item \textbf{Indices et paiements.}
\omcode{code/firm/degreeday/firm_degreeday.py}{21}{46}{Degrés-jours, paiements du swap et de l'option de vente, et burn analysis.}\label{lst:m3:freight-weather-and-other-corners:burn}
\item \textbf{La tendance.}
\omcode{code/firm/degreeday/firm_degreeday.py}{49}{56}{Correction de tendance des saisons passées vers la saison valorisée.}\label{lst:m3:freight-weather-and-other-corners:detrend}
\item \textbf{Exécuter} \texttt{m3\_weather.stats()}, \texttt{utility\_put()} et \texttt{fig\_weather.py}.
\end{tutsteps}
\textbf{Pour aller plus loin.} Valoriser la même option de vente sur les dix derniers hivers seulement, et observer de combien le prix
varie~; valoriser un swap CDD pour juillet à la même station.
\end{tutorial}

\section{Construction~: degrés-jours et règlement du fret}\label{bld:m3:freight-weather-and-other-corners:degreeday}

\begin{build}
\bfield{Rôle} La mini-entreprise vend une protection météo aux services publics et couvre son propre fret~:
elle a besoin d'indices par station, de prix de burn et du règlement des accords de fret.
\bfield{Interface} \texttt{c10\_to\_f}, \texttt{daily\_average}, \texttt{hdd}, \texttt{cdd}~;
\texttt{swap\_payoff}, \texttt{put\_payoff(index, strike, tick, cap)}~; \texttt{burn(history, payoff)}~;
\texttt{detrend(history, years, target\_year)}~; \texttt{ffa\_settlement(daily\_index, fixed\_rate, days,
lots)}.
\bfield{Règles} Températures en °F, base 65~; les saisons incomplètes sont exclues, jamais comblées~; le quantile
de burn est le quantile empirique~; le règlement d'un FFA utilise la moyenne arithmétique des jours d'indice publiés.
\bfield{Tests d'acceptation} \texttt{code/firm/degreeday/tests/}~: conversions d'unités et degrés-jours~; paiements
avec un cap~; un prix de burn sur quatre saisons~; détendance d'une série linéaire~; un règlement de FFA.
\bfield{Pour aller plus loin} Simulation de la température journalière (moyenne saisonnière plus bruit autorégressif)~; un portefeuille
de stations aux hivers corrélés~; des courbes de FFA avec saisonnalité.
\end{build}

\begin{omsources}
\item NOAA National Centers for Environmental Information, GHCN-Daily, station USW00094846 (domaine public).
\item CME Group, ``Overview of weather markets'' et chapitre 403 du règlement de marché (degree days index futures).
\item Baltic Exchange, communiqué de 2016 sur l'acquisition par SGX~; spécifications du contrat de fret de SGX.
\item S.~Jewson et A.~Brix, \emph{Weather Derivative Valuation}, Cambridge University Press, 2005.
\item CME Rulebook, chapitre 403, CME Degree Days Index Futures (copie des Internet Archives).
\end{omsources}

\section{Exercices}

\begin{exercise}[$\star$]\label{exo:m3:freight-weather-and-other-corners:1}
Une journée affiche un maximum de 40~°F et un minimum de 22~°F. Calculer ses HDD et ses CDD.
\end{exercise}

\begin{exercise}[$\star$]\label{exo:m3:freight-weather-and-other-corners:2}
Un voyage rapporte 1.5~million de dollars de fret, coûte 450\,000~\$ de carburant et 100\,000~\$ de frais de port, et dure
38 jours. Calculer son TCE.
\end{exercise}

\begin{exercise}[$\star$]\label{exo:m3:freight-weather-and-other-corners:3}
Pourquoi un FFA se règle-t-il sur une moyenne mensuelle plutôt que sur l'indice d'un seul jour~?
\end{exercise}

\begin{exercise}[$\star\star$]\label{exo:m3:freight-weather-and-other-corners:4}
Un swap HDD a un strike de 4\,700 et un tick de 20\,000~\$. L'hiver s'établit à 4\,054. Qui paie qui,
et combien~?
\end{exercise}

\begin{exercise}[$\star\star$]\label{exo:m3:freight-weather-and-other-corners:5}
Un armateur vend 30 jours d'un FFA à temps à 21\,000~\$ par jour~; la moyenne du mois s'établit à 19\,500~\$. Que
reçoit-il, et que couvre-t-il~?
\end{exercise}

\begin{exercise}[$\star\star$]\label{exo:m3:freight-weather-and-other-corners:6}
Pourquoi l'analyse de burn sur des hivers détendancés est-elle plus chère, pour cette option de vente, que sur des hivers bruts~?
\end{exercise}

\begin{exercise}[$\star\star\star$]\label{exo:m3:freight-weather-and-other-corners:7}
\emph{Programmation.} Avec \texttt{winters}, donner le nombre de saisons complètes, leur moyenne et leur écart type,
les plus chaudes et les plus froides, et la tendance ajustée par décennie.
\end{exercise}

\begin{exercise}[$\star\star\star$]\label{exo:m3:freight-weather-and-other-corners:8}
\emph{Repérer l'erreur.} «~Nous valorisons les options météo avec Black--Scholes sur l'indice HDD, à partir de sa volatilité
historique.~»
\end{exercise}

\section{Problème~: l'hiver chaud}

\begin{problem}[{Problème du week-end --- la protection hivernale d'un service public de Chicago}]\label{pb:m3:freight-weather-and-other-corners:1}
Un distributeur de gaz perd environ 20\,000~\$ de marge pour chaque HDD dont un hiver reste en deçà de la normale. Il
envisage une option de vente HDD sur O'Hare pour novembre 2026 à mars 2027, à 20\,000~\$ par HDD, plafonnée à 20~millions de dollars.

\medskip\noindent\textbf{Partie I --- L'historique.}
\begin{enumerate}
\item Combien de saisons complètes la station fournit-elle depuis 1990/91, et quelle est leur moyenne~?
\item Quel est l'écart type des totaux hivernaux~?
\item Quels hivers ont été les plus chauds et les plus froids, avec leurs totaux~?
\item Quelle tendance une régression linéaire fait-elle apparaître~?
\item Pourquoi le prix doit-il utiliser un historique détendancé~?
\end{enumerate}

\medskip\noindent\textbf{Partie II --- Le prix de burn.}
\begin{enumerate}[resume]
\item Donner l'estimation de burn détendancée de l'hiver 2026/27.
\item Donner le strike inférieur de 10~\% à cette estimation.
\item Dans quelle part des hivers ajustés l'option de vente paie-t-elle~?
\item Donner son paiement moyen (le prix de burn) et son paiement au 90\up{e}~percentile.
\item Quel a été son paiement le plus élevé, et dans quel hiver~?
\end{enumerate}

\medskip\noindent\textbf{Partie III --- La transaction.}
\begin{enumerate}[resume]
\item Pourquoi un vendeur demandera-t-il plus que le prix de burn~?
\item Que protège le cap, et combien coûte-t-il au service public~?
\item En quoi la perte du service public diffère-t-elle du paiement de l'option (base)~?
\item Un swap serait-il préférable à une option de vente pour le service public~? Pour ses actionnaires~?
\item Qui pourrait vendre l'option de vente, et comment la couvrirait-il~?
\end{enumerate}

\medskip\noindent\textbf{Partie IV --- Jugement.}
\begin{enumerate}[resume]
\item Que laissent de côté trente ans de données~?
\item Comment une prévision à long terme modifierait-elle le prix en octobre~?
\item Pourquoi les marchés météo restent-ils petits~?
\item Énoncer le \emph{résultat nommé}~: le strike de l'analyse de burn et le paiement de l'option de vente pour un événement décennal.
\item En une phrase~: qu'est-ce qu'un \omterm{def:m3:freight-weather-and-other-corners:wd}{dérivé climatique}~?
\end{enumerate}
\end{problem}

\section{Questions d'entretien}

\begin{interviewq}[$\star$ \iqroles{trader}]\label{iq:m3:freight-weather-and-other-corners:1}
Qu'est-ce qu'un degré-jour de chauffage, et qui achète une protection HDD~?
\end{interviewq}

\begin{interviewq}[$\star$ \iqroles{trader, chercheur}]\label{iq:m3:freight-weather-and-other-corners:2}
Qu'est-ce qu'un \omterm{def:m3:freight-weather-and-other-corners:ffa}{forward freight agreement}, et sur quoi se règle-t-il~?
\end{interviewq}

\begin{interviewq}[$\star\star$ \iqroles{chercheur}]\label{iq:m3:freight-weather-and-other-corners:3}
Comment valoriseriez-vous une option météo, et pourquoi pas avec Black--Scholes~?
\end{interviewq}

\begin{interviewq}[$\star\star$ \iqroles{trader}]\label{iq:m3:freight-weather-and-other-corners:4}
Un armateur disposera d'un navire dans six mois. Comment se couvre-t-il, et quels risques subsistent-ils~?
\end{interviewq}

\begin{interviewq}[$\star\star$ \iqroles{risque}]\label{iq:m3:freight-weather-and-other-corners:5}
Vous vendez des options de vente HDD sur dix villes américaines. Quel est votre risque, et comment le mesurez-vous~?
\end{interviewq}

\begin{interviewq}[$\star\star\star$ \iqroles{chercheur, développeur}]\label{iq:m3:freight-weather-and-other-corners:6}
Concevez un pipeline qui valorise quotidiennement des contrats sur degrés-jours pour n'importe quelle station, des données météo brutes jusqu'aux cotations.
\end{interviewq}
